\section{Estados financieros base}
Los ratios financieros se obtienen principalmente del Estado de Situación Financiera y del Estado de Resultados. El primero informa activos, pasivos y patrimonio en una fecha determinada. El segundo muestra ingresos, costos, gastos y utilidad de un periodo.

\begin{table}[h]
\centering
\begin{tabular}{p{0.32\textwidth}p{0.55\textwidth}}
\toprule
Estado financiero & Información que aporta \\
\midrule
Estado de Situación Financiera & Activos corrientes, activos no corrientes, pasivos corrientes, pasivos no corrientes y patrimonio. \\
Estado de Resultados & Ventas, costo de ventas, gastos, utilidad operativa, gastos financieros, impuesto y utilidad neta. \\
Estado de Flujos de Efectivo & Entradas y salidas reales de efectivo por actividades de operación, inversión y financiamiento. \\
\bottomrule
\end{tabular}
\end{table}

Para ratios que usan promedios, como cuentas por cobrar promedio, inventario promedio, activo promedio o patrimonio promedio, se requiere información inicial y final. Si solo se tiene el saldo final, el cálculo puede ser útil como aproximación didáctica, pero debe declararse la limitación.
